\documentclass[11pt]{article}
\usepackage[utf8]{inputenc}
\usepackage[margin=1in]{geometry}
\usepackage{listings}
\usepackage{xcolor}

\definecolor{codegreen}{rgb}{0,0.6,0}
\definecolor{codegray}{rgb}{0.5,0.5,0.5}
\definecolor{codepurple}{rgb}{0.58,0,0.82}

\lstset{
    commentstyle=\color{codegreen},
    keywordstyle=\color{magenta},
    numberstyle=\tiny\color{codegray},
    stringstyle=\color{codepurple},
    basicstyle=\ttfamily\small,
    breaklines=true,
    captionpos=b,
    keepspaces=true,
    numbers=left,
    numbersep=5pt,
    showspaces=false,
    showstringspaces=false,
    showtabs=false,
    tabsize=2
}

\title{The Bridge: Theory to Implementation}
\author{Connecting OOP Pillars to Functional Operations}
\date{April 29, 2026}

\begin{document}

\maketitle

\section{Overview}
Software architecture relies on the synergy between structural theory and mechanical execution. We achieve this by mapping the four core pillars of Object-Oriented Programming (OOP)—Abstraction, Encapsulation, Herencia, and Polimorfismo—onto four fundamental operations: Instantiate, Assign, Compare, and Loop.

\section{Implementation Mechanics}

\subsection{Instantiate (Creating the Object)}
Instantiating is the act of bringing a blueprint to life in memory. This is where Herencia manifests, as you can instantiate a child class that inherits all properties of its parent.

\begin{lstlisting}[language=Java]
// Instantiate: Building a concrete object from a class blueprint
Guitar myGuitar = new Guitar();
\end{lstlisting}

\subsection{Assign (Managing the Secret)}
Assignment is the mechanical tool for Encapsulation. By using methods to assign values, we keep the internal state a "secret" and ensure data integrity.

\begin{lstlisting}[language=Java]
// Assign: Using encapsulation to set a value securely
myGuitar.setVolume(11);
\end{lstlisting}

\subsection{Compare (Driving Logic)}
Comparison allows us to leverage Polimorfismo. We compare states or types to decide which specific behavior an object should execute.

\begin{lstlisting}[language=Java]
// Compare: Decision making based on object state
if (myGuitar.getVolume() > 10) {
    applyDistortion();
}
\end{lstlisting}

\subsection{Loop (Efficiency via Abstraction)}
Looping works hand-in-hand with Abstraction. We can iterate through a collection of complex objects and trigger a simple, shared command without worrying about the underlying implementation.

\begin{lstlisting}[language=Java]
// Loop: Applying abstracted methods across multiple instances
for (Guitar g : studioRack) {
    g.tune(); // Abstraction: We know it tunes, we don't care how
}
\end{lstlisting}

\section{Conclusion}
By mastering these four operations, the "secrets" of encapsulation and the logic of inheritance become functional, scalable code.

\end{document}
